\documentclass[11pt,a4paper]{amsart}
\usepackage[T1]{fontenc}
\usepackage{amssymb}
\usepackage[margin=1in]{geometry}
\usepackage[hidelinks]{hyperref}
\hypersetup{pdftitle={An acyclic obstruction to submodular threshold representation}}
\urlstyle{same}
\DeclareFontShape{OT1}{cmr}{bx}{n}{<-> cmbx10}{}
\DeclareMathSizes{8}{8}{8}{8}
\DeclareMathSizes{9}{9}{8}{8}
\DeclareMathSizes{10.95}{10.95}{8}{8}
\newtheorem{theorem}{Theorem}
\newtheorem{proposition}[theorem]{Proposition}
\newcommand{\R}{\mathbb R}
\title{An acyclic obstruction to submodular threshold representation}
\date{}
\begin{document}
\begin{abstract}
We exhibit a nine-member family in the Boolean lattice on six elements
such that each pair has a meet or join in the family and the full
dependency digraph is acyclic. Three submodularity inequalities rule
out a strict sublevel representation. The same family works in its
generated eighteen-element sublattice. This answers a question of
Elbracht, Kneip and Teegen. We also give a separate adaptation of a
related example of Simmons, Bansal and Cheriyan.
\end{abstract}
\maketitle

\section{The question}

For a subfamily $P$ of a finite lattice $L$, structural submodularity
means that for each $a,c\in P$ at least one of $a\vee c$ and
$a\wedge c$ is in $P$. We call the family order-induced if
$P=\{x\in L:f(x)<k\}$ for a submodular function $f:L\to\R$
and a threshold $k\in\R$. Submodularity means
\[
 f(a\vee c)+f(a\wedge c)\le f(a)+f(c)\qquad(a,c\in L).
\]
On a finite lattice we may add the same constant to $f$ and $k$
to make both nonnegative.

The dependency digraph of $P$ has vertex set $L$ and includes
every arc from $L\setminus P$ to $P$. For $a,b$ with the same
membership in $P$, it includes $a\to b$ exactly when some
$c\in P$ gives one of the following,
\[
 b=a\vee c,\quad a\wedge c\notin P;\qquad
 b=a\wedge c,\quad a\vee c\notin P.
\]
If $f$ induces $P$, then $f(a)>f(b)$ on every arc. This follows
directly for an outside-to-inside arc. For the other two
types, the unused corner has value at least $k>f(c)$, and
the submodular inequality gives the strict decrease.
Elbracht, Kneip and Teegen prove this in \cite[Lemma~4.1]{ekt}.
Their Question~4.5 asks whether a structurally submodular $P$
can have an acyclic dependency digraph and still fail to be
order-induced. This is also Elbracht's
Question~5.1.6 \cite[p.~199]{elbracht} and catalog Q335 \cite{catalog}.
The known distributive example of \cite[Theorem~4.3]{ekt} uses
a dependency cycle. In their Figure~2 a cycle lies entirely in
the outer digraph, while the inner digraph is acyclic.

\begin{theorem}\label{thm:main}
There is a nine-member structurally submodular family $P$ in
the Boolean lattice on six ground elements for which the entire
dependency digraph is acyclic and no real-valued submodular function
induces $P$ by a strict threshold. The same holds in the
eighteen-element sublattice generated by $P$.
\end{theorem}

\section{An explicit witness}

Encode each subset of $\{0,\ldots,5\}$ by its bitmask, so joins
are bitwise OR and meets are bitwise AND. Take
\[
 P=\{1,5,7,9,25,44,45,47,61\}.
\]
For comparable members both corners lie in $P$. The following table
contains the remaining pairs and a corner in $P$ for each.
\[
\begin{array}{rr|r@{\qquad}rr|r}
 a&c&\text{corner}&a&c&\text{corner}\\ \hline
 1&44&45&7&45&47\\
 5&9&1&7&61&5\\
 5&25&1&9&44&45\\
 5&44&45&25&44&61\\
 7&9&1&25&45&61\\
 7&25&1&25&47&9\\
 7&44&47&47&61&45
\end{array}
\]

To check acyclicity, the following table gives layers from sources
to sinks. Let $h(a)$ be the index of the row containing $a$.
\[
\begin{array}{r|l}
 0&18,22,26,30,50,54,58\\
 1&10,14,19,20,28,34,38,42,48,51,52,56,62\\
 2&2,6,16,24,32,35,49,55,59\\
 3&23,27,33,36,40,43,46,53,60\\
 4&11,21,31,39,57\\
 5&3,12,17,37,41\\
 6&4,8\\
 7&0,15,29\\
 8&13,63\\
 9&7,25,44\\
 10&47,61\\
 11&5,9\\
 12&1\\
 13&45
\end{array}
\]
Every outside point lies before every inside point, covering all
$55\cdot9=495$ outside-to-inside arcs. The remaining check runs
through $a=0,\ldots,63$ and $c\in P$. For each of the two ordered
assignments $(b,d)=(a\vee c,a\wedge c)$ or $(a\wedge c,a\vee c)$,
if $d\notin P$ and $a,b$ have equal membership in $P$, the table
satisfies $h(a)<h(b)$. This gives 294 distinct same-membership arcs.
These finite checks can be read from the bitmask table or reproduced
by the integer replay. Consequently the full digraph, including all 789
arcs, is acyclic.

Suppose that $P=\{x:f(x)<k\}$ for a submodular $f$. Apply
submodularity to the three pairs with the given corners.
\[
\begin{array}{rr|rr}
 a&c&a\vee c&a\wedge c\\ \hline
 7&25&31&1\\
 31&44&63&12\\
 12&1&13&0
\end{array}
\]
Adding the inequalities cancels $f(1)$, $f(12)$ and $f(31)$, leaving
\[
 f(0)+f(13)+f(63)\leq f(7)+f(25)+f(44).
\]
Each left term is at least $k$, and each right term is less than $k$.
This contradicts the displayed inequality.

Closing $P$ under both operations gives the sublattice
\[
 \{0,1,4,5,7,8,9,12,13,15,25,29,31,44,45,47,61,63\}.
\]
The lattice operations agree with those in the full Boolean lattice,
and the auxiliary witnesses in the arc rules are all in $P$.
Its dependency digraph is therefore the induced subgraph on the
eighteen retained masks. All terms in the three inequalities are
retained as well. This proves Theorem~\ref{thm:main} for both lattices.

\section{A related published construction}

Simmons, Bansal and Cheriyan construct pliable families and give an
example on the ground set $\{0,\ldots,7\}$ with fifteen members
\cite[Example~2]{sbc}. Their condition concerns crossing pairs; the
displayed family $F$ contains the disjoint singletons $\{3\},\{4\}$,
while it contains neither their union nor their intersection.
Their Proposition~8 rules out a symmetric submodular threshold with
a sum of difference inequalities. The theorem here concerns an
arbitrary submodular threshold and the full dependency digraph.

Their example can be adapted to give another witness to the same
lattice question. Delete the unused ground element $0$, put
$U=\{1,\ldots,7\}$, and take
\[
 P'=F\cup\{U\setminus A:A\in F\}\cup\{\varnothing,U\}.
\]
The exact replay computes all corners of the 32-member family and
all 3,494 dependency arcs on $2^U$. The arcs are acyclic and sort
into fourteen layers. If a submodular $f$ induced $P'$, then
$g(A)=f(A)+f(U\setminus A)$ would be symmetric, submodular and induce
the same $P'$ at twice the threshold.

For such a $g$, submodularity applied to $A$ and $U\setminus B$
gives $g(A\setminus B)+g(B\setminus A)\le g(A)+g(B)$.
Write $V_1=\{1,3,5,7\}$, $V_2=\{2,3,6,7\}$ and
$V_3=\{4,5,6,7\}$. Use the $k=3$ specialization of
\cite[Proposition~8]{sbc}: apply the difference inequality to
$(V_1,V_2)$, $(\{1,5\},V_3)$ and $(\{2,6\},\{4,6,7\})$.
After cancellation this gives
\[
 g(\{1\})+g(\{2\})+g(\{4,7\})\le g(V_1)+g(V_2)+g(V_3).
\]
All three right sets are in $F$. The left sets and their
complements are absent from $F$, since they are absent in the
displayed list and their complements have size at least five.
Thus they are absent from $P'$. This contradicts the strict threshold.
The adaptation uses the earlier family and inequality, together
with closure and acyclicity checks for the enlarged family.
The six-element witness above uses nine members and ordinary meet
and join inequalities. We make no priority or minimality claim.

\appendix
\section{Finite infeasibility certificates}

For a finite lattice $L$ and a nonempty proper subset $P$, let
$e_x$ be the coordinate vector for $x$ and set
\[
 D(a,b)=e_a+e_b-e_{a\vee b}-e_{a\wedge b}.
\]
If a strict threshold exists, all differences between an outside
value and an inside value are positive. The minimum of these
finitely many differences is positive, so we may rescale $f$.
Conversely, a unit gap gives a strict threshold by taking the
least outside value. Thus order inducibility is equivalent to
\begin{align*}
 D(a,b)\cdot f&\geq 0 &&(a,b\in L),\\
 (e_t-e_s)\cdot f&\geq 1 &&(s\in P,\ t\in L\setminus P).
\end{align*}
By Farkas' lemma \cite{farkas}, infeasibility is equivalent to a
nonnegative combination of these constraint rows with zero
coefficient vector and positive sum of right-hand sides. In
particular, it is equivalent to a nonnegative combination of
$D$ rows giving a nonzero vector $h$ with $h_s\geq0$ on $P$
and $h_t\leq0$ outside $P$. For the forward direction this
is the negative of the combined threshold rows. For the
reverse direction a combination producing $h$ would imply
$h\cdot f\geq0$, whereas the signs and $\sum_x h_x=0$ give
$h\cdot(f-k)<0$ under a strict threshold.

The coefficients can be chosen rational. Append the right-hand
side to each constraint row and normalize the combination to
give $(0,1)$. The augmented generators have their first
$N=|L|$ coordinates summing to zero, so lie in a subspace of
dimension $N$. Conic Carath\'eodory \cite[Section~17]{rock}
gives a representation by at most $N$ linearly independent
generators. With rational generators and rational target, the
coefficients of such a representation are rational. At least
one generator is a threshold row, so at most $N-1$ distinct
$D$ rows are needed. Clearing denominators produces an
integer certificate. This is an application of classical
finite-dimensional duality.

In a Boolean lattice, elementary square rows suffice
\cite{fujishige}. For sets $A,B$, put $C=A\cap B$ and enumerate
the elements of $A\setminus B$ and $B\setminus A$. Add the
elements of $A\setminus B$ in order to $C$. For each added
element $i$, compare its marginal at the current prefix
with its marginal after adjoining $B\setminus A$ one
element $j$ at a time. The differences telescope into
elementary square rows with lower corner omitting $i,j$.
Summing then telescopes along the $A\setminus B$ prefix to
$D(A,B)$. Each Boolean row preserves the multiplicity of each
ground element. For the witness above the three rows expand
into nine elementary squares. The paper's proof in the
eighteen-element sublattice uses the original three rows.

\newpage
\section*{Review and disclosure}
Internally reviewed, not peer reviewed. There is no proof-assistant
formalization. Under human direction, ChatGPT Astra 6.2 (OpenAI)
generated the discovery proof and original replay. OpenAI Codex
performed the independent mathematical review. GPT-6.2 Sol (OpenAI)
checked the literature and prepared this manuscript, expanded replay
and worked illustration. DannyExperiments is the project's citation
identity. The manuscript has no byline.

\begin{thebibliography}{9}
\bibitem{elbracht}
C. Elbracht, \emph{Distinguishing and witnessing dense structures in
graphs and abstract separation systems}, dissertation, Universit\"at
Hamburg, 2021.
\href{https://www.math.uni-hamburg.de/spag/dm/papers/Dissertation_Christian_Elbracht.pdf}{Author copy}.
\bibitem{ekt}
C. Elbracht, J. L. Kneip and M. Teegen, The structure of submodular
separation systems, 2021.
\href{https://arxiv.org/abs/2103.13162v2}{arXiv:2103.13162v2}.
\bibitem{catalog}
Open Problems Catalog, Q335.
\href{https://github.com/10m-u/open-problems-catalog/blob/ca8d4d44ba3db527f19f28928c240168f545ecfe/problems/mathematics/combinatorics-graph-theory-part-1.md\#q335}{Pinned entry}.
\bibitem{sbc}
M. Simmons, I. Bansal and J. Cheriyan, Symmetric submodular functions,
uncrossable functions, and structural submodularity, 2025.
\href{https://arxiv.org/abs/2601.00140v1}{arXiv:2601.00140v1}.
\bibitem{farkas}
J. Farkas, Theorie der einfachen Ungleichungen,
\emph{J. Reine Angew. Math.} 124 (1902), 1--27.
\href{https://doi.org/10.1515/crll.1902.124.1}{DOI}.
\bibitem{rock}
R. T. Rockafellar, \emph{Convex Analysis},
Princeton University Press, 1970.
\bibitem{fujishige}
S. Fujishige, \emph{Submodular Functions and Optimization}, 2nd ed.,
Annals of Discrete Mathematics 58, Elsevier, 2005.
\end{thebibliography}
\end{document}
